\section{From a Guarded Heuristic to a Natively Multi-Objective Learned Controller}
\label{sec:truesmart}

\subsection{Why the first agent needed a guard}
The first (v1) neural controller (Section~\ref{sec:mdp}) chose one of five fixed window multipliers from a five-value state, and its reward scored only the agent's own delivery rate, queueing delay and loss. Nothing in the reward mentions other flows, so the policy that is best for one flow is free to be a \emph{bandwidth bully}. We measured this (Section~\ref{sec:fairness}): against CUBIC it took $1.20\times$, $1.58\times$ and $1.73\times$ its fair share with one, two and three competitors. The fix in that version was an inference-time rule, \texttt{--polite}, that overrides the network when queueing delay rises. That rule works against CUBIC but is a hand-written patch on top of a learned policy: it is not what the network learned, it over-corrects against a single competitor, and it hides the question we care about, namely whether the learned controller can \emph{itself} be fair. This section replaces both the discrete action set and the post-hoc guard.

\subsection{State space}
Time is divided into control intervals of one smoothed round-trip time, clamped to $[100,500]$\,ms. At the end of interval $t$ the sender computes the vector $s_t\in\mathbb{R}^8$ below. Let $\mathrm{SRTT}_t$ be the smoothed RTT, $\overline{\mathrm{RTT}}_t$ the mean RTT of the interval, $\mathrm{RTT}_{\min}$ the minimum RTT seen, $qd_t=(\overline{\mathrm{RTT}}_t-\mathrm{RTT}_{\min})/\mathrm{RTT}_{\min}$ the relative queueing delay, and $w_t$ the window.
\begin{center}\footnotesize
\setlength{\tabcolsep}{3pt}
\begin{tabular}{@{}cp{2.7cm}p{4.3cm}@{}}
\toprule
 & feature & definition (clipped to $[-1,1]$ or $[0,1]$)\\
\midrule
$s^{(1)}$ & SRTT gradient & $(\mathrm{SRTT}_t-\mathrm{SRTT}_{t-1})/\mathrm{RTT}_{\min}$\\
$s^{(2)}$ & queue-delay derivative & EWMA (weight $0.5$) of $qd_t-qd_{t-1}$\\
$s^{(3)}$ & queue-delay level & $qd_t/3$\\
$s^{(4)}$ & loss ratio, last window & loss events / new packets in the interval, $/0.3$\\
$s^{(5)}$ & in-flight / BDP & in-flight $/(\max_{10}\text{rate}\cdot\mathrm{RTT}_{\min})$, $/4$\\
$s^{(6)}$ & competing-traffic pressure & Fano factor (var/mean) of ACK counts in 10\,ms bins, $/6$\\
$s^{(7)}$ & previous action & $a_{t-1}$\\
$s^{(8)}$ & window size & $\log_2 w_t/7$\\
\bottomrule
\end{tabular}
\end{center}
The pressure indicator $s^{(6)}$ is derived from ACK arrival timing: a lone paced flow receives an evenly spaced ACK stream, while cross traffic interleaves its packets with ours at the shared queue and makes our ACK stream burstier. $s^{(7)}$ and $s^{(8)}$ go beyond the five requested signals. $s^{(8)}$ is essential and not cosmetic: a policy that outputs a multiplicative factor can only realise an \emph{additive} increase, which is what makes AIMD converge to equal shares, if it knows the window. A first version without $s^{(8)}$ was trained with the same reward and failed (Section~\ref{sec:truesmart-abl}).

\subsection{Continuous action space}
The policy outputs $a_t\in(-1,1)$ and the window is updated multiplicatively,
\begin{align}
C_t&=2^{a_t}\in(0.5,\,2.0),\\
w_{t+1}&=\min\{w_{\max},\ \max\{2,\ C_t\,w_t\}\},\quad w_{\max}=128.
\end{align}
The network is a $8\!\to\!32\!\to\!32\!\to\!1$ tanh MLP (1{,}377 weights) whose output $z_t$ is the mean of a Gaussian during training, $a_t\sim\mathcal{N}(\tanh z_t,\sigma^2)$, and is used deterministically, $a_t=\tanh z_t$, at run time. The bound $C\in(0.5,2)$ per interval and the window limits are the only hard safety constraints.

\subsection{Native multi-objective reward}
For learned flow $i$ in interval $t$, with $T$ the interval length, $\mu$ the bottleneck rate and $x_j$ the delivery rate of flow $j$ among the $n$ active flows,
\begin{multline}
R_t=w_1\,\frac{x_i}{\mu}\;-\;w_2\,\min\{qd_t,2\}\;-\;w_3\,\min\{\ell_t,0.3\}\\
-\;w_4\,g_t\Bigl[(1-J_t)+\beta\max\{0,\ s_i-1\}\Bigr],
\end{multline}
\begin{align}
J_t&=\frac{(\sum_j x_j)^2}{n\sum_j x_j^2},\qquad s_i=\frac{x_i}{\frac1n\sum_j x_j},\\
g_t&=\operatorname{clip}\Bigl(\frac{\sum_j x_j/\mu-0.8}{0.15},0,1\Bigr).
\end{align}
The four terms are throughput, delay, loss and fairness penalty.
Jain's index $J_t$ is symmetric, so by itself it does not distinguish the bully from the victim. The one-sided term $\beta\max\{0,s_i-1\}$ charges the flow that is over its equal share, and $g_t$ switches the penalty off while the link is not saturated, so a lone flow, or one using capacity a competitor cannot, is never penalised for sending fast. The weights used for the reported agent are $w_1=1$, $w_2=0.2$, $w_3=3$, $w_4=1.5$, $\beta=0.5$. The fairness term uses the other flows' rates, which only the simulator knows; this is centralised training with decentralised execution, and at run time the policy sees only the local state above.

\subsection{Training}
We train with Proximal Policy Optimisation (clip $0.2$, $\gamma=0.95$, GAE $\lambda=0.95$, Adam with learning rate $10^{-3}$, six epochs per batch, entropy bonus $0.003$) with an asymmetric actor--critic: the critic additionally sees five scenario descriptors (flow count, link rate, RTT, buffer/BDP, loss) that the actor never sees. Each iteration collects 32 episodes. The environment is a packet-level, 1\,ms-tick dumbbell simulator written in C (\texttt{sim/netsim.c}) that calls the same feature code as the sender; every episode randomises the link rate ($300$--$2000$ packets/s), RTT ($20$--$150$\,ms), buffer ($0.5$--$3$ BDP), random loss ($0$--$3\%$), jitter, an optional mid-episode capacity change, one to three background flows drawn from \{CUBIC, a simplified BBR, AIMD, copies of the learned policy\}, and late-joining flows. Training ran 300 iterations (about 8 minutes on two CPU cores). The checkpoint (iteration 280) was selected on simulator scenarios with held-out seeds; all results in Section~\ref{sec:fairness} come from the real UDP stack, not the simulator.

\subsection{Deployment: weights as a C header}
\texttt{train\_smart\_poise.py --export} writes the weights to \texttt{neural\_weights.h}, a header of \texttt{static const float} arrays, which the sender includes at compile time (a text file with the same numbers can be loaded at run time instead). The forward pass is two $8\!\times\!32$ and $32\!\times\!32$ matrix--vector products and a dot product, about $1.4$k multiply--adds once per control interval, with no Python and no heap allocation. We checked that the compiled-in network and the NumPy network agree to six decimals on test inputs. In \texttt{--mode smart} there is no \texttt{--polite} rule and no hand-over to AIMD: after slow start ends the policy alone sets the window. What is \emph{not} learned is listed in Section~\ref{sec:truesmart-lim}.

\subsection{Ablations}
\label{sec:truesmart-abl}
\textbf{Without the window in the state} (7 inputs, otherwise identical, three reward settings and two critic designs): the agent drifted to a small window and was starved, taking roughly $0.2\times$ its fair share on held-out simulator scenarios; the logs are kept in \texttt{results/smart/ablation\_nocwnd\_*.csv}. \textbf{Without the fairness term} ($w_4=0$, everything else identical, same 300 iterations): the agent learned to be a bully in the simulator, taking $1.8\times$, $2.7\times$ and $3.6\times$ its share against one, two and three CUBIC flows with Jain's index falling to $0.30$. This ablation is the evidence that the fairness reward, and not the architecture, removes the bully behaviour.

\subsection{What is still not learned}
\label{sec:truesmart-lim}
Slow start (classic $+1$ per ACK until the window reaches $64$ or the RTT ratio reaches $1.3$ or a loss occurs), pacing at $2w/\mathrm{SRTT}$, the retransmission timer and fast retransmit, and the window bounds $[2,128]$ are fixed rules shared with the baselines. The reward weights and the held-out-seed checkpoint selection were chosen by us after a small number of trials (about ten training runs of 200--300 iterations in total, including the failed ones). The simulator's BBR and CUBIC are our own simplified implementations.
